\clearpage
\fancyhead[L]{\small CARLINI. MAILÄNDER EPHEMERIDEN FÜR 1816.}
\SourceLeaf{597}
\newcommand{\ObservedLibelleDenominator}{\hbox{\vrule width0pt height4.600000bp depth0pt\special{pdf:bcontent}\special{pdf:code q 0 g 2.146667 4.600000 m 2.146667 4.293333 l 2.453333 3.986667 l 2.453333 3.680000 l 1.840000 3.066667 l 1.533333 3.066667 l 0.613333 2.146667 l 0.000000 2.146667 l 0.000000 0.613333 l 0.306667 0.613333 l 0.613333 0.306667 l 1.226667 0.306667 l 1.533333 0.000000 l 3.066667 0.000000 l 3.986667 0.920000 l 3.986667 1.840000 l 3.066667 2.760000 l 3.066667 3.066667 l 3.680000 3.680000 l 3.680000 3.986667 l 3.066667 4.600000 l h 1.533333 2.146667 m 1.840000 2.453333 l 2.453333 1.840000 l 2.453333 1.533333 l 1.840000 0.920000 l 1.533333 1.226667 l h f* Q}\special{pdf:econtent}\kern3.986667bp}}

\begin{center}\small Göttingische gelehrte Anzeigen. Stück 47. Seite 457..46c. 1816 März 23.\\[3pt]\rule{.72\linewidth}{.3pt}\end{center}

\emph{Effemeridi astronomiche di Milano per l’anno bisestile} 1816, \emph{calcolate da} \textsc{Francesco Carlini}. \emph{Con appendice. Dalla cesarea regia stamperia.} Mailand 1815. Octav. Die Ephemeriden 108 Seiten, der Anhang 100 Seiten.

Die Zusätze, welche diesmal den astronomischen Kalender begleiten, sind 1) das Verzeichniss der 34 Fundamentalsterne nebst ihren jährlichen Präcessionen und den Constanten für die Berechnung ihrer Aberration und Nutation, wie im Jahrgange für 1815 (m. s. unsere Anzeigen vom v. J. S. 89 [S. 574 d. B.]). 2) Die \textsc{Carlini}'sche Refractionstafel. Diese Tafel, welche auf die \textsc{Laplace}'sche Hypothese über den Zustand der Atmosphäre und die daraus abgeleitete Theorie der Refraction gegründet ist, jedoch nach einer neuen auf eignen Beobachtungen beruhenden Bestimmung der dabei vorkommenden Constanten, war schon im Jahrgange der Ephemeriden für 1808 abgedruckt. Da Rec. diesen nicht zur Hand hat, so kann er nicht angeben, ob die gegenwärtige Tafel ein blos unveränderter Abdruck ist. 3) Verzeichniss von Fixsternbedeckungen im Jahre 1816 für den Meridian und die Polhöhe von Florenz von den Florenzer Astronomen berechnet.

Der Anhang enthält folgende Aufsätze. 1) Beobachtete Refraction bei kleinen Höhen über dem Horizont, von \textsc{Barnabas Oriani}. Unter den in den vorhergehenden Jahrgängen bekannt gemachten Beobachtungen mit dem dreifussigen \textsc{Reichenbach}'schen Multiplicationskreise befinden sich 19 im Jahre 1811 beobachtete untere Culminationen der Capella, welcher Stern in Mailand in einer Höhe von $1^\circ\;36'$, nur wenige Minuten über den den Horizont begrenzenden Bergen durchgeht. Ueberdies waren im Jahr 1811 noch 14 obere Culminationen beobachtet, woraus mit der Polhöhe $45^\circ\;28'\;0''7$ die mittlere Declination $45^\circ\;47'\;28''5$ für den Anfang des Jahrs 1811 abgeleitet wurde. Die untern Culminationen wurden von Hrn. \textsc{Oriani} nach den vornehmsten Refractionstafeln berechnet, nemlich den Tafeln von \textsc{Bradley}, \textsc{Tob. Mayer}, \textsc{Piazzi}, \textsc{Delambre}, \textsc{Carlini} und \textsc{Bessel}, unter denen die von \textsc{Carlini} am besten und bis auf ein paar Secunden übereinstimmten, in sofern nemlich die vermeinte Verbesserung für die nordliche Hälfte des Meridians, die \textsc{Carlini} früherhin für nöthig gehalten hatte, weggelassen wurde. 2) Fortsetzung der Beobachtungen über das periodische Schwanken der Gebäude, von \textsc{Angelo Cesaris}. Die schon im Jahrgange der Ephemeriden für 1813 mitgetheilten Erfahrnngen über diesen Gegenstand (m. s. unsre Anzeigen 1813 S. 951 [S. 556 d. B.]) werden hier mit andern nicht weniger merkwürdigen vermehrt. Dort war von einer Bewegung die Rede, die gleichsam wie eine Rotation der Mauern, an welchen Mauerquadrant und Mitagsfernrohr aufgehängt sind, um eine verticale Axe angesehen werden können: die neuen Erfahrungen beziehen sich auf eine Rotation um eine horizontale von Osten nach Westen gerichtete Axe. Diese Bewegung wurde sichtbar gemacht durch eine vortreffliche \textsc{Reichenbach}sche am Mauerquadranten angebrachte Libelle von einer solchen Empfindlichkeit, dass eine Veränderung von Einer Secunde einen Ausschlag von $1\frac{2}{\ObservedLibelleDenominator}$ Linien gab. Dieselben Veränderungen, welche sich an dieser Libelle zeigten, liessen sich auch an der mit dem Mauerquadranten beobachteten Zenithdistanz, so wie auch an einer zweiten an der entgegengesetzten Seite der Mauer angebrachten Libelle erkennen. Die hier mitgetheilten Erfahrungen deuten sehr bestimmt auf einen Zusammenhang mit dem Wetter und mit der Stärke der Erwärmung des Gebäudes durch das Sonnenlicht. Am regelmässigsten sind die Veränderungen von Vormittags bis Nachmittags. Bei heiterm Wetter bewegt sich die Blase in der Libelle nach Süden zu; die Grösse der Bewegung ist im Durchschnitt etwa $2''$, und bei bedecktem Himmel fällt sie ganz weg. 3) Tafel für die Mittelpunktsgleichung der Vesta (Excentricität $=0.0889$) und \SourceLeaf{598}für die Reduction auf die Ekliptik (Neigung der Bahn $=7^\circ\;8'\;20''$, von \textsc{Carlini}. 4) Sternbedeckungen, beobachtet zu Mailand in den Jahren 1811--1815 von demselben, und zu Florenz in den Jahren 1810--1813 von den dortigen Astronomen. --- 5) Schiefe der Ekliptik aus den mit dem \textsc{Reichenbach}'schen dreifussigen Kreise angestellten Solstitialbeobachtungen, von \textsc{Barnabas Oriani}. Die Beobachtungen im Wintersolstitium 1810 und in den beiden Solstitien von 1811 waren bereits in den frühern Jahrgängen mitgetheilt: hier erhalten wir die ausführlichen Beobachtungen in den fünf folgenden Solstitien, und die aus allen gezogenen Resultate. Es ist sehr interessant, nun auch den Ausspruch der grössern \textsc{Reichenbach}'schen Vervielfältigungskreise über diesen vielbesprochenen Gegenstand zu erhalten. Die vier Wintersolstitien geben die mittlere Schiefe der Ekliptik auf den Anfang von 1812 reducirt $=23^\circ\;27'\;48''20$, die vier Sommersolstitien $23^\circ\;27'\;50''77$, also noch immer einen obwohl viel kleinern Unterschied, als die zwölfzolligen \textsc{Reichenbach}'schen Kreise und andere Instrumente gegeben haben. Zur Reduction der Beobachtungen wurden \textsc{Carlini}'s Refractionstafeln gebraucht: mit den \textsc{Delambre}'schen Refractionstafeln würde der Unterschied etwas grösser sein und $3''71$ betragen. --- 6) Meteorologische Beobachtungen im Jahr 1814, von \textsc{Angelo Cesaris}, nebst einer Uebersicht des gefallenen Regens in den Jahren 1764--1814. Die mittlere Quantität des gefallenen Regens aus allen 51 Jahren findet sich 35 Zoll 3,92 Linien: die Quantität in den einzelnen Jahren läuft natürlich sehr unordentlich; nimmt man sie aber für eine längere Reihe von Jahren zusammen, so zeigt sich im Ganzen eine stete Zunahme. \textsc{Cesaris} schreibt dies der immer mehr verbreiteten Bewässerung der Felder zu, wodurch eine grosse Quantität Wasser über eine viel grössere Oberfläche verbreitet einer stärkern Verdunstung ausgesetzt werde und so wieder in häufigern Regen zurückkommen müsse. Die Anzahl der Jahre ist indess wohl noch zu klein, und der Erfolg vom Wechsel des Zufalls zu abhängig, um sichere allgemeine Resultate ziehen zu können. Ganz ausgezeichnet feucht war das letzte Jahr 1814, wo die Menge des gefallenen Regens 58 Zoll 11,58 Linien betrug; die kleinste Quantität im Jahre 1771 war 25 Zoll 11,6 Linien.
